% BAB I - TENTANG APLIKASI
\section{Tentang Aplikasi}

\subsection{Deskripsi Proyek}

Proyek ini mengembangkan sistem multimedia yang memungkinkan pengguna menghasilkan musik hanya dengan gerakan tangan secara real-time. Sistem ini memanfaatkan MediaPipe Hand Tracking untuk mendeteksi posisi dan gerakan tangan, yang kemudian digunakan untuk mengontrol dua elemen utama: Arpeggiator dan Drum Machine.

Pada Arpeggiator, tinggi posisi tangan mengatur pitch nada, sementara gerakan pinch gesture mengatur volume suara. Sedangkan Drum Machine memungkinkan pengguna untuk memainkan instrumen drum secara individual dengan mengangkat jari yang sesuai.

\subsection{Anggota Tim}

Proyek \textbf{Gestune: Musik dari Gerakan Tangan} ini dikembangkan oleh tim yang terdiri dari:

\begin{table}[H]
\centering
\caption{Anggota Tim Pengembang}
\label{tab:team}
\begin{tabular}{lll}
\hline
\textbf{Nama Lengkap} & \textbf{NIM} & \textbf{ID GitHub} \\ \hline
A Edwin Krisandika Putra & 122140003 & \href{https://github.com/aloisiusedwin}{aloisiusedwin} \\
Fathan Andi Kartagama & 122140055 & \href{https://github.com/pataanggs}{pataanggs} \\
Rizki Alfariz Ramadhan & 122140061 & \href{https://github.com/Alfariz11}{Alfariz11} \\
\hline
\end{tabular}
\end{table}

\subsection{Fitur Utama}

Aplikasi Gestune memiliki beberapa fitur utama yang memungkinkan pengguna untuk membuat musik dengan gerakan tangan:

\subsubsection{Arpeggiator (Tangan Kiri)}
\begin{itemize}
    \item \textbf{Kontrol Pitch}: Tinggi posisi tangan kiri mengatur pitch nada yang dihasilkan. Semakin tinggi posisi tangan, semakin tinggi pitch nada yang dihasilkan.
    \item \textbf{Kontrol Volume}: Gesture pinch (mendekatkan ibu jari dan telunjuk) mengatur volume suara. Semakin kecil jarak antara ibu jari dan telunjuk, semakin kecil volume suara, begitu juga sebaliknya.
\end{itemize}

\subsubsection{Fungsi Gesture Tangan Kanan}
\begin{itemize}
    \item \textbf{Pemicu bunyi instrumen drum}: Setiap jari tangan kanan merepresentasikan satu instrumen drum sehingga perubahan posisi jari dapat langsung menghasilkan bunyi \textit{kick}, \textit{snare}, \textit{hihat}, \textit{hightom}, atau \textit{crash cymbal}. Ketika jari berada pada posisi terbuka, instrumen terkait akan mengikuti pola ritmis yang telah diprogram pada \textit{drum sequencer}.
    \item \textbf{Pergantian pola ritme}: Gesture menggenggam (\textit{fist gesture}) digunakan untuk mengganti pola ritme drum secara otomatis tanpa menggunakan tombol fisik atau antarmuka tambahan.
    \item \textbf{Kontrol BPM}: Fitur ini bekerja dengan membuka kunci pengatur BPM. Saat tangan kanan membuat gesture pinch, pengguna dapat menggerakkan tangan ke atas atau ke bawah untuk menyesuaikan BPM (Beats Per Minute) dari pola ritme drum yang sedang dimainkan.
\end{itemize}